\section{Conclusion}

\paragraph{A}
This paper has reconstructed capacity utilization as a structurally mediated object under capitalist disequilibrium. Departing from treatments that view utilization either as an exogenous engineering constant, an unmediated behavioral ratio, or an empirical regression residual, we have demonstrated that capacity utilization ($\mu_t = Y_t / Y_t^p$) becomes identifiable only downstream of the recovery of the aggregate transformation elasticity $\theta_t$.

\paragraph{B}
By integrating Okishio’s theory of decentralized market anarchy, Vidal’s framework of the internally divided firm, and Marglin-style distributive grounding, we established that the transformation of capital accumulation into potential productive capacity ($g_{Y^p} = \theta_t g_K$) is endogenously conditioned by class distribution, corporate financialization, capital composition, and external balance constraints.

\paragraph{C}
Empirically, our cointegrating polynomial regression (CMPR IM-OLS) estimations delivered two major historical conclusions:

\begin{enumerate}
  \item \textbf{Center (United States, 1929–2024):} The U.S. corporate economy operates permanently within a structural regime of capital overaccumulation ($\theta_t < 1.00$). The historical transition from post-war Fordism to neoliberal post-Fordism represents an endogenous regime shift—driven by wage suppression, relative machinery price inflation, and corporate financialization—that generated secular downward drift in capacity utilization.
\end{enumerate}

\begin{enumerate}
  \item \textbf{Periphery (Chile, 1920–2024):} In the peripheral economy, the transformation relation is non-linearly gated by foreign exchange availability. Using Chen’s (2015) Profile Least Squares algorithm over the Prebischian Capacity to Import ($C_t = X_t \cdot ToT_t$), we demonstrated that when external constraints bind ($D^{ext}_t = 1$), the Okishio distributive dividend breaks down, triggering a severe peripheral realization trap ($\hat{\delta}_2 = -15.99^{***}$).
\end{enumerate}

\paragraph{D}
These findings demonstrate that the difference between center and periphery is not a difference between a realized norm and a structural deviation, but a difference in how the same fundamental transformation problem is mediated under distinct world-system conditions.

\paragraph{E}
The scope of this chapter is deliberately bounded. It does not present a complete crisis theory, a fully specified model of profitability dynamics, or an exhaustive dependency framework. The chapter reconstructs capacity utilization as a structurally mediated object and identifies the transformation elasticity as the hinge of its empirical recovery across center and periphery, while leaving the broader crisis and profitability architecture for subsequent work.